% Source: https://artofproblemsolving.com/community/c3h2423270_crazy_challenging_problem
% Stable identity: aops-topic-2423270-q01
Let $r$ be an integer root, so $c=-\left(ar^2+br\right)=-r(ar+b)$. If $r=0$ then $c=0$ and $a^2+b^2=0$, so $a=0$, which is excluded. So $r\ne0$. Put $u=ar+b$, so $b=u-ar$ and $c=-ru$. The condition $a^2+b^2=c^2$ becomes
\[a^2+(u-ar)^2=r^2u^2\iff\left(1+r^2\right)a^2-2ru\,a+\left(1-r^2\right)u^2=0.\]
If $u=0$, this gives $\left(1+r^2\right)a^2=0$, so $a=0$, excluded. So $u\ne0$, and the quadratic equation in $a$ has the rational root $a$, so its discriminant
\[4r^2u^2-4\left(1+r^2\right)\left(1-r^2\right)u^2=4u^2\left(r^4+r^2-1\right)\]
must be a perfect square, i.e. $r^4+r^2-1$ must be a perfect square. For $|r|\ge2$,
\[\left(r^2\right)^2<r^4+r^2-1<\left(r^2+1\right)^2,\]
so it lies strictly between consecutive squares and is not a square. Hence $r=\pm1$.
\begin{itemize}
\item $r=1$: the equation is $2a^2-2ua=0$, so $a=u$ (as $a\ne0$), $b=u-a=0$ and $c=-u=-a$.
\item $r=-1$: the equation is $2a^2+2ua=0$, so $a=-u$, $b=u+a=0$ and $c=u=-a$.
\end{itemize}
In both cases $(a,b,c)=(a,0,-a)$. Conversely, for every nonzero integer $a$, the triple $(a,0,-a)$ works: $a^2+0^2=(-a)^2$, and $ax^2-a=0$ has the integer solutions $x=\pm1$. So the answer is
\[\boxed{(a,b,c)=(a,0,-a),\ a\ne0}.\]
Since $b=0$ in every solution, no triple of positive integers works.
